\section{Требования к программе}

\subsection{Функциональные требования}

\headingtextsep

Программная система должна обеспечивать следующие функциональные возможности.

\subsubsection{Управление рабочими пространствами и участниками}

Система должна обеспечивать создание и настройку рабочих пространств.

Система должна обеспечивать управление составом участников, ролями и правами доступа в пределах рабочего пространства.

Система должна обеспечивать раздельное администрирование на уровне рабочего пространства и на уровне платформы.

\subsubsection{Управление задачами и зависимостями}

Система должна обеспечивать создание, редактирование и архивирование задач.

Система должна поддерживать статусы, приоритеты, сроки, назначение исполнителя и группировку задач.

Система должна обеспечивать хранение зависимостей между задачами и поддержку представлений списка, доски, иерархии и других форм отображения на основе одной модели данных.

\subsubsection{Документация, заметки и файлы}

Система должна обеспечивать ведение проектной документации в формате markdown-страниц.

Система должна обеспечивать хранение заметок, вложенных файлов и изображений с возможностью связывания их с задачами и документами.

Система должна обеспечивать навигацию по иерархии документов и переходы по внутренним ссылкам.

\subsubsection{Поиск, история и интеграции}

Система должна обеспечивать полнотекстовый и семантический поиск по задачам, документам, заметкам и связанным сущностям.

Система должна обеспечивать ведение журнала действий и истории изменений.

Система должна обеспечивать интеграционное взаимодействие с внешними сервисами и агентами через API и MCP-совместимый контур.

\subsubsection{Режимы функционирования}

Основной режим:

\begin{itemize}
    \item выполнение полного набора пользовательских и агентских операций;
    \item сохранение изменений, файлов и событий аудита;
    \item работа интеграций и механизмов поиска.
\end{itemize}

Локальный режим разработки:

\begin{itemize}
    \item запуск системы с упрощенной схемой аутентификации для отладки;
    \item использование локального хранилища файлов;
    \item возможность запуска без внешних облачных учетных данных.
\end{itemize}

\textheadingsep
\subsection{Требования к интерфейсу}

\headingtextsep

\subsubsection{Пользовательский веб-интерфейс}

\paragraph{Адаптивность}

Веб-интерфейс должен корректно работать на экранах настольных и мобильных устройств в пределах поддерживаемых разрешений.

\paragraph{Локализация}

На первом этапе система должна поддерживать русский язык интерфейса.

\paragraph{Общие принципы оформления}

Интерфейс должен быть строгим, функциональным и ориентированным на работу с насыщенными информацией экранами.

\subsubsection{Экран аутентификации}

Система должна предоставлять отдельный экран входа для пользователей с отображением понятных сообщений об ошибках аутентификации.

\subsubsection{Главная рабочая область}

После входа пользователь должен получать доступ к основным разделам рабочего пространства: задачам, документации, файлам, заметкам, поиску и журналу активности.

\subsubsection{Интерфейс работы с задачами}

Интерфейс задач должен поддерживать фильтрацию, сортировку, поиск, изменение статуса, просмотр зависимостей и переход к связанным сущностям.

\subsubsection{Интерфейс базы знаний и файлов}

Интерфейс документации должен поддерживать иерархию страниц, просмотр markdown-содержимого, работу с вложениями и переходы по связанным материалам.

\subsubsection{Интерфейс администрирования рабочего пространства}

Интерфейс администрирования должен обеспечивать управление участниками, ролями, интеграциями, агентами и учетными данными агентов.

\subsubsection{Навигация и уведомления}

Система должна использовать единые навигационные элементы, хлебные крошки, уведомления об успешных операциях и понятные сообщения об ошибках.

\textheadingsep
\subsection{Требования к реализации}

\headingtextsep

Система должна быть реализована по клиент-серверной архитектуре.

Серверная часть должна быть реализована на языке Rust.

Пользовательский веб-интерфейс должен быть реализован с использованием React и TypeScript.

Основное хранение структурированных данных должно выполняться в PostgreSQL.

Хранение файлов должно выполняться через абстракцию, поддерживающую локальный backend и S3-совместимое объектное хранилище.

Развертывание системы должно поддерживаться в контейнеризированной среде.

\textheadingsep
\subsection{Требования к надежности}

\headingtextsep

Система должна обеспечивать сохранность пользовательских данных при штатном завершении работы и при перезапуске сервисов.

Система должна обеспечивать журналирование ошибок и событий аудита, достаточное для диагностики отказов.

Система должна поддерживать резервное копирование базы данных и файловых объектов.

При временной недоступности внешних интеграций система должна сохранять работоспособность основных внутренних функций.

\textheadingsep
\subsection{Требования к тестированию}

\headingtextsep

\subsubsection{Модульное тестирование}

Для серверной и клиентской частей должны быть предусмотрены модульные тесты ключевой бизнес-логики.

\subsubsection{Интеграционное тестирование}

Должны быть предусмотрены интеграционные тесты сценариев работы API, аутентификации, поиска и хранения данных.

\subsubsection{Проверка пользовательских сценариев}

Должны быть предусмотрены сценарии проверки основных пользовательских потоков веб-интерфейса.

\textheadingsep
\subsection{Требования к установке}

\headingtextsep

Система должна поставляться в виде набора конфигурационных файлов и инструкций, достаточных для локального и серверного развертывания.

Установка должна поддерживать запуск в локальной среде разработки и в контейнеризированной среде.

\textheadingsep
\subsection{Требования к сопровождению}

\headingtextsep

Система должна сопровождаться эксплуатационной документацией, журналами изменений и инструкциями по обновлению.

Должны быть предусмотрены процедуры миграции схемы данных и восстановления системы после сбоя.

\textheadingsep
\subsection{Требования к безопасности}

\headingtextsep

Система должна обеспечивать аутентификацию пользователей через внешний провайдер идентификации и поддерживать локальный режим разработки с отдельными ограничениями.

Учетные данные агентов должны храниться в виде защищенных хешей и поддерживать отзыв, ротацию и ограничение областей доступа.

Веб-интерфейс должен использовать защищенные сессии, защиту от CSRF-атак и аудит критичных административных операций.

\textheadingsep
\subsection{Требования к документации}

\headingtextsep

По проекту должны быть разработаны и актуализироваться следующие документы:

\begin{itemize}
    \item техническое задание;
    \item описание архитектуры системы;
    \item руководство пользователя;
    \item руководство разработчика или программиста;
    \item инструкция по развертыванию и администрированию.
\end{itemize}

Документация должна вестись на русском языке и храниться в формате, пригодном для дальнейшего сопровождения и обновления.